% ==============================================================================
% preamble_template.tex - תבנית קדם-מסמך מאומתת לספר אקדמי / דוח מכין ב-LuaLaTeX
% מבוסס על כללי "מהבעיה אל ה-PDF" ושיעור L03 (ד"ר יורם סגל)
% ==============================================================================

\usepackage{geometry}
\geometry{
  a4paper,
  top=2.5cm,
  bottom=2.5cm,
  inner=3.0cm,   % 2.5cm + 0.5cm Gutter (שוליים + תוספת כריכה לפי טבלה 11.1)
  outer=2.5cm,
  headheight=16pt
}

% מתמטיקה ונוסחאות
\usepackage{amsmath,amssymb,amsfonts}

% כיווניות ושפות ב-LuaLaTeX באמצעות babel + bidi=basic (תקן L03 ופרק 12)
\usepackage[bidi=basic,provide=*,shorthands=off]{babel}
\babelprovide[import,main]{hebrew}
\babelprovide[import]{english}
\usepackage{fontspec}

% הגדרת גופנים תקניים מתוך מערכת ההפעלה Windows (לפי טבלה 11.1 וטבלה 12.1)
\babelfont{rm}[
  Path=C:/Windows/Fonts/,
  UprightFont=times.ttf,
  BoldFont=timesbd.ttf,
  ItalicFont=timesi.ttf,
  BoldItalicFont=timesbi.ttf
]{Times New Roman}

\babelfont[hebrew]{rm}[
  Script=Hebrew,
  Path=C:/Windows/Fonts/,
  UprightFont=david.ttf,
  BoldFont=davidbd.ttf,
  ItalicFont=david.ttf,
  BoldItalicFont=davidbd.ttf
]{David}

\babelfont{sf}[
  Path=C:/Windows/Fonts/,
  UprightFont=arial.ttf,
  BoldFont=arialbd.ttf,
  ItalicFont=ariali.ttf,
  BoldItalicFont=arialbi.ttf
]{Arial}

\babelfont[hebrew]{sf}[
  Script=Hebrew,
  Path=C:/Windows/Fonts/,
  UprightFont=arial.ttf,
  BoldFont=arialbd.ttf,
  ItalicFont=ariali.ttf,
  BoldItalicFont=arialbi.ttf
]{Arial}

\babelfont{tt}[
  Path=C:/Windows/Fonts/,
  UprightFont=cour.ttf,
  BoldFont=courbd.ttf
]{Courier New}

\babelfont[hebrew]{tt}[
  Script=Hebrew,
  Path=C:/Windows/Fonts/,
  UprightFont=cour.ttf,
  BoldFont=courbd.ttf
]{Courier New}

\newcommand{\hebrewfont}{\rmfamily}
\newcommand{\hebrewfontsf}{\sffamily}
\newcommand{\hebrewfonttt}{\ttfamily}
\newcommand{\englishfont}{\rmfamily}
\newcommand{\englishfontsf}{\sffamily}
\newcommand{\englishfonttt}{\ttfamily}

% תאימות מלאה לפקודות polyglossia/babel עבור טקסט דו-כיווני (BiDi) - בטוח גם במצב מתמטיקה
\providecommand{\texthebrew}[1]{\ifmmode\mbox{\foreignlanguage{hebrew}{#1}}\else\foreignlanguage{hebrew}{#1}\fi}
\providecommand{\textenglish}[1]{\ifmmode\mathrm{#1}\else\foreignlanguage{english}{#1}\fi}
\newenvironment{hebrew}{\begin{otherlanguage}{hebrew}}{\end{otherlanguage}}
\newenvironment{english}{\begin{otherlanguage}{english}}{\end{otherlanguage}}

% ריווח שורות ומניעת כותרות יתומות (טבלה 11.1 ופרק 19.4)
\linespread{1.28}\selectfont
\newcommand{\onehalfspacing}{\linespread{1.28}\selectfont}
\newcommand{\needspace}[1]{\par\nopagebreak\vspace{0pt}}

% תמיכה בפרמטרים אופציונליים ברשימות enumerate ו-itemize
\let\origenumerate\enumerate
\let\endorigenumerate\endenumerate
\renewenvironment{enumerate}[1][]{\origenumerate}{\endorigenumerate}
\let\origitemize\itemize
\let\endorigitemize\enditemize
\renewenvironment{itemize}[1][]{\origitemize}{\endorigitemize}

% צבעים ארגוניים עקביים לכל המסמך (פרק 15 ופרק 16.6)
\usepackage{xcolor}
\definecolor{BrandNavy}{RGB}{18, 48, 89}
\definecolor{BrandBlue}{RGB}{32, 92, 168}
\definecolor{BrandLightBlue}{RGB}{235, 243, 252}
\definecolor{BrandTeal}{RGB}{16, 112, 104}
\definecolor{BrandLightTeal}{RGB}{234, 247, 245}
\definecolor{BrandGold}{RGB}{184, 115, 20}
\definecolor{BrandLightGold}{RGB}{253, 247, 234}
\definecolor{BrandRed}{RGB}{178, 34, 45}
\definecolor{BrandLightRed}{RGB}{253, 237, 238}
\definecolor{BrandGray}{RGB}{244, 246, 249}
\definecolor{TableHeaderBg}{RGB}{218, 229, 244}

% גרפיקה, איורים, TikZ ו-pgfplots (פרק 16)
\usepackage{graphicx}
\graphicspath{{figures/}}
\usepackage{tikz}
\usetikzlibrary{arrows.meta,positioning,shapes.geometric,calc,patterns,backgrounds,fit,babel}
\usepackage{pgfplots}
\pgfplotsset{compat=1.18}

% טבלאות מעוצבות (פרק 11 ופרק 16.5)
\usepackage{array}
\usepackage{tabularx}
\newcommand{\toprule}{\noalign{\hrule height 1.1pt\vspace{3pt}}}
\newcommand{\midrule}{\noalign{\vspace{2pt}\hrule height 0.6pt\vspace{3pt}}}
\newcommand{\bottomrule}{\noalign{\vspace{2pt}\hrule height 1.1pt}}
\newcommand{\addlinespace}[1][4pt]{}
\newcommand{\rowcolors}[3]{}
\newcommand{\rowcolor}[1]{}
\newcolumntype{R}[1]{>{\raggedleft\arraybackslash}p{#1}}
\newcolumntype{L}[1]{>{\raggedright\arraybackslash}p{#1}}
\newcolumntype{C}[1]{>{\centering\arraybackslash}p{#1}}
\newcolumntype{Y}{>{\raggedleft\arraybackslash}X}

% הגנה מלאה על מספור פרקים, סעיפים, איורים, טבלאות ומשוואות מפני היפוך BiDi (טבלה 14.3 וטבלה 15.5)
\setcounter{secnumdepth}{3}
\setcounter{tocdepth}{3}
\renewcommand{\thechapter}{\textenglish{\arabic{chapter}}}
\renewcommand{\thesection}{\textenglish{\arabic{chapter}.\arabic{section}}}
\renewcommand{\thesubsection}{\textenglish{\arabic{chapter}.\arabic{section}.\arabic{subsection}}}
\renewcommand{\thesubsubsection}{\textenglish{\arabic{chapter}.\arabic{section}.\arabic{subsection}.\arabic{subsubsection}}}
\renewcommand{\thefigure}{\textenglish{\arabic{chapter}.\arabic{figure}}}
\renewcommand{\thetable}{\textenglish{\arabic{chapter}.\arabic{table}}}
\renewcommand{\theequation}{\arabic{chapter}.\arabic{equation}}

% פקודת כותרת רצה הניתנת להתאמה בכל פרויקט
\providecommand{\ProjectRunningTitle}{\texthebrew{דוח מכין לפרויקט מחקר אקדמי}}
\providecommand{\ProjectCourseName}{\texthebrew{סמינריון מחקרי בניהול טכנולוגיות ומימון}}

% עיצוב כותרות פרקים וסעיפים וכותרות רצות (ללא תלות בחבילות חיצוניות)
\makeatletter
\def\@makechapterhead#1{%
  \vspace*{10\p@}%
  {\parindent \z@ \raggedleft \normalfont\sffamily
    \ifnum \c@secnumdepth >\m@ne
      \Large\bfseries\color{BrandBlue} \@chapapp\space \thechapter
      \par\nobreak
      \vskip 6\p@
    \fi
    \interlinepenalty\@M
    {\color{BrandNavy}\hrule height 1.5pt}\vskip 8\p@
    \Huge \bfseries \color{BrandNavy} #1\par\nobreak
    \vskip 8\p@
    {\color{BrandNavy}\hrule height 0.8pt}
    \vskip 22\p@
  }}
\def\@makeschapterhead#1{%
  \vspace*{10\p@}%
  {\parindent \z@ \raggedleft \normalfont\sffamily
    \interlinepenalty\@M
    {\color{BrandNavy}\hrule height 1.5pt}\vskip 8\p@
    \Huge \bfseries \color{BrandNavy} #1\par\nobreak
    \vskip 8\p@
    {\color{BrandNavy}\hrule height 0.8pt}
    \vskip 22\p@
  }}
\renewcommand\section{\@startsection {section}{1}{\z@}%
  {-3.2ex \@plus -1ex \@minus -.2ex}%
  {1.8ex \@plus.2ex}%
  {\normalfont\sffamily\Large\bfseries\color{BrandNavy}}}
\renewcommand\subsection{\@startsection{subsection}{2}{\z@}%
  {-2.8ex\@plus -1ex \@minus -.2ex}%
  {1.2ex \@plus .2ex}%
  {\normalfont\sffamily\large\bfseries\color{BrandBlue}}}
\renewcommand\subsubsection{\@startsection{subsubsection}{3}{\z@}%
  {-2.2ex\@plus -1ex \@minus -.2ex}%
  {0.8ex \@plus .2ex}%
  {\normalfont\sffamily\normalsize\bfseries\color{BrandTeal}}}

% כותרות עליונות ותחתונות עם כיוון מפורש בכל משבצת (טבלה 14.3 ופרק 15.10)
\def\ps@fancy{%
  \def\@oddhead{\vbox to \headheight{\hbox to \textwidth{\footnotesize\color{BrandNavy}\ProjectRunningTitle\hfill\color{BrandBlue}\ProjectCourseName}\vss\hrule height 0.5pt}}%
  \def\@evenhead{\@oddhead}%
  \def\@oddfoot{\vbox{\hrule height 0.4pt\vskip 3pt\hbox to \textwidth{\footnotesize\color{BrandNavy}\textenglish{\copyright\ 2026}\ \texthebrew{כל הזכויות שמורות לצוות המחקר}\hfill\texthebrew{עמוד }\textenglish{\thepage}}}}%
  \def\@evenfoot{\@oddfoot}%
}
\def\ps@plain{%
  \def\@oddhead{}%
  \def\@evenhead{}%
  \def\@oddfoot{\vbox{\hrule height 0.4pt\vskip 3pt\hbox to \textwidth{\footnotesize\color{BrandNavy}\textenglish{\copyright\ 2026}\ \texthebrew{כל הזכויות שמורות לצוות המחקר}\hfill\texthebrew{עמוד }\textenglish{\thepage}}}}%
  \def\@evenfoot{\@oddfoot}%
}
\makeatother
\pagestyle{fancy}

% תיבות מעוצבות (tcolorbox) - מסגרת ב-LTR ותוכן ב-RTL לפי כלל 6 בפרומפט B1 (פרק 14.7)
\usepackage{tcolorbox}

% 1. תיבת "למה זה חשוב למנהל / לקובע המדיניות"
\newenvironment{managerbox}[1][למה זה חשוב למנהל ולקובע המדיניות]{%
  \par\medskip%
  \begin{english}%
  \begin{tcolorbox}[
    colback=BrandLightGold,
    colframe=BrandGold,
    colbacktitle=BrandGold,
    coltitle=white,
    boxrule=1pt,
    arc=3pt,
    left=10pt,right=10pt,top=8pt,bottom=8pt,
    title={\hfill\textbf{\sffamily\texthebrew{#1}}}
  ]%
  \begin{hebrew}%
}{%
  \end{hebrew}%
  \end{tcolorbox}%
  \end{english}%
  \par\medskip%
}

% 2. תיבת הגדרה / מודל מתמטי / מסר מרכזי
\newenvironment{modelbox}[1][הגדרה ומודל אנליטי]{%
  \par\medskip%
  \begin{english}%
  \begin{tcolorbox}[
    colback=BrandLightBlue,
    colframe=BrandNavy,
    colbacktitle=BrandNavy,
    coltitle=white,
    boxrule=1pt,
    arc=3pt,
    left=10pt,right=10pt,top=8pt,bottom=8pt,
    title={\hfill\textbf{\sffamily\texthebrew{#1}}}
  ]%
  \begin{hebrew}%
}{%
  \end{hebrew}%
  \end{tcolorbox}%
  \end{english}%
  \par\medskip%
}

% 3. תיבת מלכודת מתודולוגית / אזהרת סיכון
\newenvironment{warningbox}[1][מלכודת מתודולוגית ומוקד סיכון]{%
  \par\medskip%
  \begin{english}%
  \begin{tcolorbox}[
    colback=BrandLightRed,
    colframe=BrandRed,
    colbacktitle=BrandRed,
    coltitle=white,
    boxrule=1pt,
    arc=3pt,
    left=10pt,right=10pt,top=8pt,bottom=8pt,
    title={\hfill\textbf{\sffamily\texthebrew{#1}}}
  ]%
  \begin{hebrew}%
}{%
  \end{hebrew}%
  \end{tcolorbox}%
  \end{english}%
  \par\medskip%
}

% 4. תיבת יישום בפרויקט המחקר ("בפרויקט שלנו")
\newenvironment{projectbox}[1][יישום בפרויקט המחקר שלנו]{%
  \par\medskip%
  \begin{english}%
  \begin{tcolorbox}[
    colback=BrandLightTeal,
    colframe=BrandTeal,
    colbacktitle=BrandTeal,
    coltitle=white,
    boxrule=1pt,
    arc=3pt,
    left=10pt,right=10pt,top=8pt,bottom=8pt,
    title={\hfill\textbf{\sffamily\texthebrew{#1}}}
  ]%
  \begin{hebrew}%
}{%
  \end{hebrew}%
  \end{tcolorbox}%
  \end{english}%
  \par\medskip%
}

% פקודות עזר דו-כיווניות (BiDi) לפי Prompt B1 (פרק 14.7) - בטוחות במצב טקסט ובמצב מתמטיקה
\newcommand{\en}[1]{\ifmmode\mathrm{#1}\else\foreignlanguage{english}{#1}\fi}
\newcommand{\num}[1]{\ifmmode#1\else\foreignlanguage{english}{#1}\fi}
\newcommand{\pct}[1]{\ifmmode#1\%\else\foreignlanguage{english}{#1\%}\fi}
\newcommand{\nis}[1]{\ifmmode#1\;\mbox{\foreignlanguage{hebrew}{ש"ח}}\else\foreignlanguage{english}{#1}\ \foreignlanguage{hebrew}{ש"ח}\fi}
\newcommand{\usd}[1]{\ifmmode\$#1\else\foreignlanguage{english}{\$#1}\fi}
\newcommand{\hebmath}[1]{\mbox{\foreignlanguage{hebrew}{#1}}}

% הדפסת ביבליוגרפיה בתקן IEEE מתוך קובץ references.bib המסונכרן ל-bibliography_items.tex
\newcommand{\printbibliography}{%
  \chapter{רשימת מקורות וביבליוגרפיה (\en{References})}%
  % Auto-generated IEEE bibliography items from references.bib
\begin{english}
\begin{thebibliography}{99}
\small
\bibitem{segal2026prep} \textbf{Segal, Yoram}, ``From Problem to PDF: The Professional Preparatory Report (מהבעיה אל ה-PDF: המדריך לדוח המכין המקצועי),'' \textit{Bar-Ilan University, Department of Management}, 2026.
\bibitem{ulrich2020product} \textbf{Ulrich, Karl T. and Eppinger, Steven D. and Yang, Maria C.}, ``Product Design and Development,'' \textit{McGraw-Hill Education}, 2020.
\bibitem{boi2024annual} \textbf{Bank of Israel}, ``Bank of Israel Annual Report for 2024 -- Chapter 9: The Housing Market and Construction Sector,'' \textit{Bank of Israel Research Department}, 2025.
\bibitem{boi2025stability} \textbf{Bank of Israel}, ``Financial Stability Report -- First Half of 2025: Credit Risks in Real Estate and Developer Promotions,'' \textit{Bank of Israel Financial Stability Division}, 2025.
\bibitem{mof2025economist} \textbf{Chief Economist Department}, ``Real Estate Market Review: Analysis of Contractor Cash-Flow Promotions (20/80) and Investor Purchases,'' \textit{Israel Ministry of Finance}, 2025.
\bibitem{cbs2026dwellings} \textbf{Central Bureau of Statistics (CBS)}, ``New Dwellings Sold and Unsold Inventory in Israel: Statistical Bulletin Q1--Q3 2026,'' \textit{Israel Central Bureau of Statistics}, 2026.
\bibitem{boi2025directive} \textbf{Banking Supervision Department}, ``Temporary Order Amending Proper Conduct of Banking Business Directives 329 and 311: Restrictions on Bullet/Balloon Housing Loans and Project Accompaniment Risk Weights (March 2025),'' \textit{Bank of Israel}, 2025.
\bibitem{mian2009consequences} \textbf{Mian, Atif and Sufi, Amir}, ``The Consequences of Mortgage Credit Expansion: Evidence from the U.S. Mortgage Default Crisis,'' \textit{The Quarterly Journal of Economics}, vol.~124, no.~4, pp.~1449--1496, 2009. DOI: \texttt{10.1162/qjec.2009.124.4.1449}.
\bibitem{mayer2009rise} \textbf{Mayer, Christopher and Pence, Karen and Sherlund, Shane M.}, ``The Rise in Mortgage Defaults,'' \textit{Journal of Economic Perspectives}, vol.~23, no.~1, pp.~27--50, 2009. DOI: \texttt{10.1257/jep.23.1.27}.
\bibitem{segal2026l03} \textbf{Segal, Yoram}, ``L03: Create PDF Documents with LaTeX Vibe Coding (שיעור 3: מסמך PDF מקצועי בלטך),'' \textit{University of Haifa, Course Lecture Notes}, 2026.
\bibitem{doran1981smart} \textbf{Doran, George T.}, ``There's a S.M.A.R.T. Way to Write Management's Goals and Objectives,'' \textit{Management Review}, vol.~70, no.~11, pp.~35--36, 1981.
\bibitem{ita2025carman} \textbf{Israel Tax Authority}, ``CARMAN Real Estate Transactions Database : Methodological Documentation and Reporting Standards,'' \textit{Ministry of Finance, Jerusalem}, 2025.
\bibitem{webster2002analyzing} \textbf{Webster, Jane and Watson, Richard T.}, ``Analyzing the Past to Prepare for the Future: Writing a Literature Review,'' \textit{MIS Quarterly}, vol.~26, no.~2, pp.~xiii--xxiii, 2002.
\bibitem{stiglitz1981credit} \textbf{Stiglitz, Joseph E. and Weiss, Andrew}, ``Credit Rationing in Markets with Imperfect Information,'' \textit{The American Economic Review}, vol.~71, no.~3, pp.~393--410, 1981.
\bibitem{jensen1976theory} \textbf{Jensen, Michael C. and Meckling, William H.}, ``Theory of the Firm: Managerial Behavior, Agency Costs and Ownership Structure,'' \textit{Journal of Financial Economics}, vol.~3, no.~4, pp.~305--360, 1976. DOI: \texttt{10.1016/0304-405X(76)90026-X}.
\bibitem{minsky1992financial} \textbf{Minsky, Hyman P.}, ``The Financial Instability Hypothesis,'' \textit{The Jerome Levy Economics Institute Working Paper}, no.~74, pp.~1--10, 1992.
\bibitem{kindleberger2011manias} \textbf{Kindleberger, Charles P. and Aliber, Robert Z.}, ``Manias, Panics, and Crashes: A History of Financial Crises,'' \textit{Palgrave Macmillan}, 2011.
\bibitem{reinhart2009this} \textbf{Reinhart, Carmen M. and Rogoff, Kenneth S.}, ``This Time Is Different: Eight Centuries of Financial Folly,'' \textit{Princeton University Press}, 2009.
\bibitem{keys2010did} \textbf{Keys, Benjamin J. and Mukherjee, Tanmoy and Seru, Amit and Vig, Vikrant}, ``Did Securitization Lead to Lax Screening? Evidence from Subprime Loans,'' \textit{The Quarterly Journal of Economics}, vol.~125, no.~1, pp.~307--362, 2010. DOI: \texttt{10.1162/qjec.2010.125.1.307}.
\bibitem{gorton2009securitized} \textbf{Gorton, Gary B.}, ``Slapped in the Face by the Invisible Hand: Banking and the Panic of 2007,'' \textit{Federal Reserve Bank of Atlanta Financial Markets Conference}, pp.~1--43, 2009.
\bibitem{ben2011inflated} \textbf{Ben-David, Itzhak}, ``Financial Constraints and Inflated Home Prices during the Real Estate Boom,'' \textit{American Economic Journal: Applied Economics}, vol.~3, no.~3, pp.~55--87, 2011. DOI: \texttt{10.1257/app.3.3.55}.
\bibitem{ghent2011recourse} \textbf{Ghent, Andra C. and Kudlyak, Marianna}, ``Recourse and Residential Mortgage Default: Evidence from US States,'' \textit{The Review of Financial Studies}, vol.~24, no.~9, pp.~3139--3186, 2011. DOI: \texttt{10.1093/rfs/hhr055}.
\bibitem{salelaw1974} \textbf{Knesset of Israel}, ``Sale (Apartments) (Assurance of Investments of Persons Acquiring Apartments) Law, 5735-1974 (חוק המכר  , התשל"ה-1974),'' \textit{Book of Laws (Sefer HaChukim)}, 1974.
\bibitem{sandberg2011ways} \textbf{Sandberg, J\"orgen and Alvesson, Mats}, ``Ways of Constructing Research Questions: Gap-Spotting or Problematization?,'' \textit{Organization}, vol.~18, no.~1, pp.~23--44, 2011.
\bibitem{bernanke2010monetary} \textbf{Bernanke, Ben S.}, ``Monetary Policy and the Housing Bubble,'' \textit{Speech at the Annual Meeting of the American Economic Association}, pp.~1--22, 2010.
\bibitem{foote2008subprime} \textbf{Foote, Christopher L. and Gerardi, Kristopher and Goette, Lorenz and Willen, Paul S.}, ``Subprime Facts: What (We Think) We Know about the Subprime Crisis and What We Don't,'' \textit{Federal Reserve Bank of Boston Public Policy Discussion Papers}, no.~08-2, pp.~1--64, 2008.
\bibitem{pugh1991total} \textbf{Pugh, Stuart}, ``Total Design: Integrated Methods for Successful Product Engineering,'' \textit{Addison-Wesley}, 1991.
\bibitem{pmbok2021} \textbf{Project Management Institute}, ``A Guide to the Project Management Body of Knowledge (PMBOK Guide),'' \textit{Project Management Institute}, 2021.
\bibitem{iso31000} \textbf{International Organization for Standardization}, ``ISO 31000:2018 -- Risk Management -- Guidelines,'' \textit{ISO}, 2018.
\end{thebibliography}
\end{english}%
}

% קישורים פנימיים וחיצוניים (hyperref נטען אחרון)
\usepackage[unicode]{hyperref}
\hypersetup{
  colorlinks=true,
  linkcolor=BrandNavy,
  citecolor=BrandBlue,
  urlcolor=BrandTeal
}
